\documentclass[12pt]{article}
\usepackage{graphicx,import}
\usepackage[svgnames]{xcolor} 
\usepackage{fancyhdr}
\usepackage{subfig}
\usepackage{hyperref}
\usepackage{enumitem}
\usepackage{cite}
\usepackage[many]{tcolorbox}
\usepackage{listings }
\usepackage[a4paper, total={6in, 8in} , bottom = 25mm , top = 25mm, headheight = 1.25cm , includehead,includefoot,heightrounded ]{geometry}
\usepackage{afterpage}
\usepackage{amssymb}
\usepackage{pdflscape}
\usepackage{gensymb}
\usepackage{textcomp}
\usepackage{tikz,pgfplots}
\usepackage{xecolor}
\usepackage{rotating}
\usepackage{pdfpages}
\usepackage[Kashida]{xepersian}
\usepackage[T1]{fontenc}
\usepackage{tikz}
\usepackage[utf8]{inputenc}
\usepackage{PTSerif} 
\usepackage{seqsplit}

\usepackage[edges]{forest}

\usepackage{listings}
\usepackage{xcolor}

\hypersetup{
	colorlinks   = true, %Colours links instead of ugly boxes
	urlcolor     = blue, %Colour for external hyperlinks
	linkcolor    = blue, %Colour of internal links
	citecolor   = red %Colour of citations
}
 
\definecolor{codegreen}{rgb}{0,0.6,0}
\definecolor{codegray}{rgb}{0.5,0.5,0.5}
\definecolor{codepurple}{rgb}{0.58,0,0.82}
\definecolor{backcolour}{rgb}{0.95,0.95,0.92}
 
\NewDocumentCommand{\codeword}{v}{
\texttt{\textcolor{blue}{#1}}
}
\lstset{language=java,keywordstyle={\bfseries \color{blue}}}

\lstdefinestyle{mystyle}{
    backgroundcolor=\color{backcolour},   
    commentstyle=\color{codegreen},
    keywordstyle=\color{magenta},
    numberstyle=\tiny\color{codegray},
    stringstyle=\color{codepurple},
    basicstyle=\ttfamily\normalsize,
    breakatwhitespace=false,         
    breaklines=true,                 
    captionpos=b,                    
    keepspaces=true,                 
    numbers=left,                    
    numbersep=5pt,                  
    showspaces=false,                
    showstringspaces=false,
    showtabs=false,                  
    tabsize=2
}

\lstset{style=mystyle}

\settextfont[Scale=1.2 ,BoldFont={Bahij Nazanin-Bold.ttf} , ItalicFont = {IRNazaninIranic.ttf}]{Bahij Nazanin-Regular.ttf}
\setlatintextfont[Scale = 1.0]{Garamond}
\DefaultMathsDigits 
\DeclareMathSizes{11}{19}{13}{9} 
%\DeclareMathSizes{12}{14.4}{8}{9}





\newenvironment{changemargin}[2]{%
\begin{list}{}{%
\setlength{\topsep}{0pt}%
\setlength{\leftmargin}{#1}%
\setlength{\rightmargin}{#2}%
\setlength{\listparindent}{\parindent}%
\setlength{\itemindent}{\parindent}%
\setlength{\parsep}{\parskip}%
}%
\item[]}{\end{list}}


\definecolor{foldercolor}{RGB}{124,166,198}

\tikzset{pics/folder/.style={code={%
    \node[inner sep=0pt, minimum size=#1](-foldericon){};
    \node[folder style, inner sep=0pt, minimum width=0.3*#1, minimum height=0.6*#1, above right, xshift=0.05*#1] at (-foldericon.west){};
    \node[folder style, inner sep=0pt, minimum size=#1] at (-foldericon.center){};}
    },
    pics/folder/.default={20pt},
    folder style/.style={draw=foldercolor!80!black,top color=foldercolor!40,bottom color=foldercolor}
}

\forestset{is file/.style={edge path'/.expanded={%
        ([xshift=\forestregister{folder indent}]!u.parent anchor) |- (.child anchor)},
        inner sep=1pt},
    this folder size/.style={edge path'/.expanded={%
        ([xshift=\forestregister{folder indent}]!u.parent anchor) |- (.child anchor) pic[solid]{folder=#1}}, inner xsep=0.6*#1},
    folder tree indent/.style={before computing xy={l=#1}},
    folder icons/.style={folder, this folder size=#1, folder tree indent=3*#1},
    folder icons/.default={12pt},
}

\begin{document}


%%% title pages
\begin{titlepage}
\begin{center}
        
\vspace*{0.7cm}

\includegraphics[width=0.4\textwidth]{sharif1.png}\\
\vspace{0.5cm}
\textbf{ \Huge{\emph ‌سیگنال‌ها و سیستم‌ها} }\\
\vspace{0.5cm}
\textbf{ \Large{ تمرین سوم} }
\vspace{0.2cm}
       
 
      \large \textbf{دانشکده مهندسی کامپیوتر}\\\vspace{0.2cm}
    \large   دانشگاه صنعتی شریف\\\vspace{0.2cm}
       \large   ﻧﯿﻢ سال دوم 00-99 \\\vspace{0.2cm}
      \noindent\rule[1ex]{\linewidth}{1pt}
استاد:\\
    \textbf{{جناب آقای دکتر منظوری شلمانی}}


    \vspace{0.15cm}
نام و نام خانوادگی:\\

       
    \textbf{{امیرمهدی نامجو - 97107212}}
\end{center}
\end{titlepage}
%%% title pages


%%% header of pages
\newpage
\pagestyle{fancy}
\fancyhf{}
\fancyfoot{}
\cfoot{\thepage}
\chead{تمرین سوم}
\rhead{\includegraphics[width=0.1\textwidth]{sharif.png}}
\lhead{امیرمهدی نامجو}
%%% header of pages

\KashidaOff

\section{سوال اول}



$$x[n] =  u[n-1] + u[n-2] + ... u[n-N]$$

$$H(z) = \frac{z^{-1}}{1 - z^{-1}} + \frac{z^{-2}}{1 - z^{-1}} + \frac{z^{-3}}{1 - z^{-1}} + \cdots + \frac{z^{-N}}{1 - z^{-1}}$$

$$= \frac{z^{-N}}{1 -z^{-1}}(1+z+z^2 + \cdots + z^{N-1}) = \frac{z^{-N}}{1 -z^{-1}} (\frac{1-z^N}{1-z})$$

$$=\frac{z-z^{1-N}}{(z-1)^2}$$

\section{سوال دوم}


\section{سوال سوم}

$$
H(z)=\frac{1-z^{-1}}{1+\frac{3}{4} z^{-1}}
$$

\begin{enumerate}[label = \Alph*)]
	\item 
	$$H(z) = \frac{-4}{3} + \frac{\frac{7}{3}}{1 + \frac{3}{4} z^{-1}}$$
	
	$$\mathcal{Z}^{-1}(H(z))(n) = h[n] = \frac{-4}{3} \delta[n] + \frac{7}{3} (\frac{-3}{4})^n u[n]$$ 
	
	\item
	
	$$x[n] = (1/3)^n u[n] + u[-n -1]$$
	
	$$X(z) = \frac{1}{1 - (1/3) z^{-1}} - \frac{1}{1 - z^{-1}}; ROC: 1/3<|z|<1$$
	
	$$Y(z) = X(z) H(z) = (\frac{1}{1 - (1/3) z^{-1}} - \frac{1}{1 - z^{-1}}) (\frac{-4}{3} + \frac{\frac{7}{3}}{1 + \frac{3}{4} z^{-1}})$$$$ = \frac{\frac{-8}{13}}{1 - 1/3 z^{-1}} + \frac{\frac{8}{13}}{1 + 3/4 z^{-1}} ; ROC: 1/3<|z| < 3/4$$
	
	ساده سازی عبارت بالا به کمک Mathematica انجام شده است.
	
	$$y[n] = \frac{-8}{13} (1/3)^n u[n] - \frac{8}{13} (-3/4)^{n} u[-1-n] $$
	
	
	\item
	سیستم LTI علی پایدار است اگر و تنها اگر قطب هایش از نظر اندازه کمتر مساوی $1$ باشند.
	
	در این مورد $|z| = 3/4$ است که کمتر از یک است. پس پایدار است.
	
	
	
\end{enumerate}



\section{سوال چهارم}
$$
\begin{array}{l}
	x[n]=\left(\frac{1}{3}\right)^{n} u[n]+2^{n} u[-n-1] \\
	y[n]=5\left(\frac{1}{3}\right)^{n} u[n]-5\left(\frac{2}{3}\right)^{n} u[n]
\end{array}
$$



\begin{enumerate}[label = \Alph*)]
	\item 

$$X(z) = \frac{1}{1-1/3 z^{-1}} - \frac{1}{1- 2 z^{-1}} ; ROC: 1/3<|z|<2$$
$$Y(z)= \frac{5}{1 - 1/3 z^{-1}} - \frac{5}{1 - 2/3 z^{-1}} ; ROC: 1/3 < |z| <2/3$$

$$H(z) = \frac{Y(z)}{X(z)} =\frac{\frac{5}{1-\frac{1}{3 z}}-\frac{5}{1-\frac{2}{3 z}}}{\frac{1}{1-\frac{1}{3 z}}-\frac{1}{1-\frac{2}{z}}} =\frac{3 (z-2)}{3 z-2}$$
$$=3 - \frac{2}{1 - 2/3 z^{-1}} ; ROC: |z| > 2/3  $$


قطب سیستم $a = 2/3$ و صفر آن $z=2$ است.


\item
$$h[n] = 3\delta[n] - 2 (2/3)^n u[n]$$


\item

$$3z Y(z) -2 Y(z) = 3(zX(z) - 2 X(z))) \rightarrow$$
$$3 y[n+1] - 2y[n] = 3(x[n+1] - 2x[n])$$

\item

بله چون ROC تابع $H(z)$ آن شامل دایره واحد می شود.
\end{enumerate}

\section{سوال چهارم}

\end{document}



